\note{Five or six sentences, in the order of the paper: (1) a diffusion model needs a denoiser; Markov sequences
(time series, 1-D fields) fix its form. (2) Gaussian chain: diffusion breaks locality, the score is not local.
(3) Chain + noise = tree, so BP gives the exact score from the kernel. (4) Learn the kernel by EM with BP as E-step,
from noisy data. (5) Laplace chain: learned kernel vs networks (best network 2.1--2.8$\times$, MLP 5.8--9.9$\times$);
how it degrades off-Markov. (6) Allen--Cahn: the family matters; parameters recovered.}

Diffusion models learn to denoise a process at every time step $t$. To accomplish this, the score is the fundamental object that is needed, which is an affine function of the posterior mean. We ask how to do this for Markov sequences. Starting with the Gaussian case, where the score is linear and depends on the precision matrix, we see how noise quickly erases the Markov structure. In the Laplace case, the score does not have a closed form anymore, and we notice that the factor graph is a tree; hence, we can get the exact score via Belief Propagation. Messages are probability densities; we then use a grid to discretize them, and the discretization error is tested with the Gaussian closed form. Subsequently, we then want to estimate the kernel, either with Expectation Maximisation over latent variables or Maximum Likelihood over the clean sequence. Moreover, we run experiments on a suitable physical model, such as the stochastic Allen-Cahn field, comparing many kernels from different parametric families with more flexible neural kernels. Finally, we show that in Markov sequences, it can be beneficial to train small neural kernels and recover the score through BP rather than using neural networks to denoise the score directly.
